\documentclass[11pt]{article}

\usepackage[a4paper,margin=1in]{geometry}
\usepackage[T1]{fontenc}
\usepackage[utf8]{inputenc}
\usepackage{lmodern}
\usepackage{microtype}
\usepackage{amsmath}
\usepackage{booktabs}
\usepackage{enumitem}
\usepackage{hyperref}
\usepackage{xcolor}

\hypersetup{
  colorlinks=true,
  linkcolor=blue!50!black,
  urlcolor=blue!50!black,
  citecolor=blue!50!black
}

\title{Fast Geometry and Staged Transport for Large LXe TPC Simulation}
\author{LXeMC Design Note}
\date{\today}

\begin{document}
\maketitle

\begin{abstract}
This note documents the current design of the LXeMC detector and transport
model, using the LZ detector as a concrete example of a large liquid-xenon
time projection chamber (LXe TPC). The design goal is not maximal generality.
The goal is to simulate very large event samples quickly while retaining the
physics that matters for rejection and classification. The resulting approach is
a staged simulation: a specialized fast kernel performs most of the transport
outside the fiducial volume (FV), where most events are rejected quickly, while
a full photon/electron stack is retained only inside the FV, where detailed
single-site versus multi-site behavior matters.
\end{abstract}

\section{Introduction}

Large LXe TPCs are designed to reject backgrounds aggressively. In a detector
such as LZ, most external gammas are rejected by active veto regions before they
reach the fiducial volume. This observation drives the central design choice of
the present simulation: the transport outside the FV should be optimized for
speed and early rejection, while the transport inside the FV should preserve the
detailed cascade structure needed for event topology studies.

The code therefore does not attempt to use one uniform transport model
everywhere. Instead, it exploits the detector semantics directly. Outside the
FV, the simulation only needs enough geometry and physics to decide whether a
gamma is vetoed, absorbed in passive material, or survives into the FV. Inside
the FV, the simulation keeps the full stack of photons and charged particles
because the event classification depends on the spatial pattern of energy
depositions.

\section{Problem Statement}

The practical task is to simulate a very large number of gamma events in a large
LXe detector. In the intended use case, the number of events is so large that a
uniform full simulation everywhere in the detector is too expensive. This is not
primarily a software-engineering inconvenience; it is a physics-rate problem.
The simulation must be fast enough to process event counts large enough to
support background studies at the level of interest.

The immediate consequence is that the dominant cost must be spent only on the
small subset of trajectories that can still affect the final physics observable.
Any cost spent on fully tracking obviously rejected events is wasted.

\section{Constraints and Boundary Conditions}

The design is shaped by the following constraints.

\begin{enumerate}[label=\arabic*.]
\item \textbf{Speed is dominant.} The simulation must run at the level of
millions of events per second on representative benchmark configurations.

\item \textbf{Most events are rejected early.} For many source configurations,
gammas are vetoed in active xenon outside the FV, absorbed in passive material,
or degraded to energies irrelevant for the region of interest.

\item \textbf{Detailed transport is only needed locally.} The expensive part of
the physics---full photon/electron cascades and cluster structure---matters
primarily inside the FV.

\item \textbf{The detector partition is physically meaningful.} The geometry is
not arbitrary. It should reflect the detector regions that carry distinct
transport or classification meaning.
\end{enumerate}

\section{Design Philosophy}

The design is based on a simple principle: use detector semantics to reduce the
cost of transport.

The detector is partitioned into regions that matter for either material
transport or event classification. The transport kernel is then compiled
\emph{ad hoc} from that partition. This is intentionally specialized. The
detector geometry is not interpreted as a fully generic tree at runtime in the
hot path. Instead, the canonical geometry is loaded once and transformed into a
compact fast-kernel representation with explicit scalar region parameters.

This leads to three consequences.

\begin{enumerate}[label=\arabic*.]
\item The geometry used by the fast kernel is chosen for speed rather than for
maximal abstraction.

\item The detector partition is explicit and canonical. It is not inferred
dynamically from overlapping semantic flags at runtime.

\item The full transport stack is retained only where the physics demands it,
namely inside the FV.
\end{enumerate}

\section{Simulation Strategy}

The current simulation strategy has three stages.

\subsection{Fast rejection}

The first goal is to reject events as soon as they become irrelevant. If a gamma
deposits visible energy above the veto threshold in the active xenon outside the
FV, the event is rejected immediately. If a gamma degrades below the configured
ROI floor in passive material, it is terminated locally because it can no longer
produce an FV event in the energy region of interest.

\subsection{Fast transport outside the FV}

Outside the FV, a specialized fast kernel handles gamma transport using a small
set of canonical detector regions. This kernel performs:

\begin{itemize}
\item point classification into canonical detector regions,
\item analytic distance-to-boundary calculations,
\item material-aware interaction sampling,
\item immediate veto and ROI-floor decisions.
\end{itemize}

This stage is optimized for the common case in which the gamma is rejected or
absorbed before it can matter to the final physics analysis.

\subsection{Full transport inside the FV}

If a gamma survives into the FV, the fast kernel hands off its state to a
full photon/electron transport stack restricted to the FV geometry. This stack
retains:

\begin{itemize}
\item photon propagation,
\item lepton propagation,
\item secondaries through the particle stack,
\item deposit clustering relevant to event topology.
\end{itemize}

Inside the FV, the geometry is simple: a homogeneous LXe cylinder bounded by the
fiducial radial and longitudinal cuts. The expensive part is therefore not the
geometric navigation but the cascade physics itself.

\section{Current Status and Benchmark}

The current code implements the staged architecture described above.

\begin{itemize}
\item The canonical detector geometry is defined explicitly in
\texttt{detector\_lz\_v3.json}.
\item The fast kernel uses the canonical detector partition outside the FV.
\item The FV transport uses a dedicated \texttt{FVGeometry} plus a full
particle stack.
\item The legacy flat-detector workflow has been removed from the canonical
code path.
\end{itemize}

Representative benchmarks place the simulation in the target regime of roughly
$10^6$ events per second for the fast calibration workflow, with exact values
depending on source position and direction. Configurations aimed directly toward
the FV are naturally slower than configurations rejected promptly in outer veto
regions, but the code remains in the intended high-throughput regime.

\section{Geometry Model}

\subsection{Shared geometric primitives}

The file \texttt{geometry\_core.jl} contains the shared shape and wrapper layer.
It provides:

\begin{itemize}
\item primitive solids:
  \texttt{Cyl}, \texttt{CylShell}, \texttt{Box}, \texttt{Disk},
  \texttt{Cap}, \texttt{DomedContainer}, \texttt{CappedCylinder};
\item logical wrappers:
  \texttt{LCyl}, \texttt{LCylShell}, \texttt{LBox}, \texttt{LDisk};
\item materialized wrappers:
  \texttt{PhysicalVolume}, \texttt{PCyl}, \texttt{PCylShell},
  \texttt{PBox}, \texttt{PDisk};
\item common geometric operations such as containment tests, masses, and
ray entry/exit distances.
\end{itemize}

These are general geometry tools. They are not tied to the fast kernel or to the
FV stack specifically.

\subsection{Canonical detector geometry}

The file \texttt{geometry.jl} now contains the canonical detector geometry
representation. It loads the tracking detector from JSON into a tree of
\texttt{TrackingDetector} nodes, with explicit region semantics and geometry.

The canonical LZ tracking geometry is partitioned into the following regions:

\begin{itemize}
\item \texttt{AirDome}
\item \texttt{LXe\_passive}
\item \texttt{TopActive}
\item \texttt{BarrelActive}
\item \texttt{BottomActive}
\item \texttt{FV}
\item \texttt{Skin}
\item \texttt{FC\_PTFE}
\item \texttt{FC\_rings}
\end{itemize}

This partition is exhaustive for the intended detector model and reflects the
physics logic directly.

\subsection{Meaning of the canonical regions}

\begin{description}[leftmargin=2.0cm,style=nextline]
\item[\texttt{AirDome}] Gas region above the liquid xenon.
\item[\texttt{LXe\_passive}] Passive liquid xenon outside the active veto
regions and below the cathode.
\item[\texttt{TopActive}] Active xenon cylinder between the gate and the top of
the FV.
\item[\texttt{BarrelActive}] Active xenon shell around the FV over the FV
longitudinal span.
\item[\texttt{BottomActive}] Active xenon cylinder between the bottom of the FV
and the cathode.
\item[\texttt{FV}] Fiducial xenon volume used for event selection and topology.
\item[\texttt{Skin}] Active xenon skin outside the field-cage region.
\item[\texttt{FC\_PTFE}] PTFE part of the field cage.
\item[\texttt{FC\_rings}] Titanium ring part of the field cage.
\end{description}

\subsection{Source geometry}

The source geometry is described in a separate JSON file from the detector
tracking geometry. The separation is by \emph{contents}, not by dimensions:

\begin{itemize}
\item Source geometry hosts radioactive emission locations and source-side
      transport. It includes the cryostat Ti shells, MLI insulation, and the
      PMT structures whose effect is absorbed into the source-flux step before
      tracking begins.
\item Tracking geometry contains only what affects per-step photon transport
      after the source-flux step: the detector envelope, sensitive LXe regions,
      veto regions, and the field-cage structures that obstruct photons in
      transit.
\end{itemize}

When a volume appears in both files (\texttt{FC\_PTFE}, \texttt{FC\_rings},
\texttt{Skin}), it carries identical geometric dimensions on the two sides:
the source side is anchored to the bb0nu measured mass, and the tracking side
inherits the same R, z, and wall thickness. This means the fast-kernel
point-in-volume tests use anatomically correct dimensions and the
geometric-mass invariant is validated against the same reference on both
sides. Strict source/tracking dimensional consistency is enforced by the test
suite.

The split is therefore a \emph{subset} relation: tracking geometry is a subset
of source geometry, restricted to volumes the per-step transport actually
needs to see. It is not a license to draw the same volume with different
dimensions on the two sides.

\section{Transport Model}

\subsection{Fast kernel outside the FV}

The fast kernel operates on a compiled geometry representation generated from the
canonical detector tree. It performs:

\begin{enumerate}[label=\arabic*.]
\item region classification,
\item analytic distance-to-boundary calculation,
\item material lookup,
\item interaction-distance sampling,
\item immediate region-dependent decisions.
\end{enumerate}

This kernel treats transport in the outer detector as a problem of fast
geometric and material decisions, not as a full generic cascade simulation.

\subsection{Full stack inside the FV}

The FV transport uses a dedicated \texttt{FVGeometry} object compiled directly
from the canonical detector geometry. The FV stack keeps:

\begin{itemize}
\item photons,
\item electrons and positrons,
\item bremsstrahlung,
\item the particle stack and deposit list.
\end{itemize}

The geometry inside the FV is reduced to what is physically needed:
the LXe material and the cylindrical FV boundary.

\subsection{Handoff logic}

The fast kernel tracks a gamma until one of the following occurs:

\begin{itemize}
\item it is vetoed outside the FV,
\item it is absorbed or terminated below the ROI floor in passive material,
\item it escapes,
\item it reaches the FV and remains relevant.
\end{itemize}

Only in the last case is the gamma handed to the FV transport stack.

\subsection{ROI floor}

The simulation uses an ROI floor in gamma energy. Once a gamma energy drops
below this threshold, the code assumes that it can no longer generate a relevant
contribution to the final ROI and terminates it locally. This is a deliberate
analysis-motivated approximation used to gain speed.

\subsection{Veto logic}

The fast kernel uses region-dependent visible-energy thresholds:

\begin{itemize}
\item active xenon outside the FV:
  veto at the configured TPC threshold;
\item skin:
  veto at the configured skin threshold;
\item passive regions:
  no veto, only transport and possible local termination.
\end{itemize}

\section{Event Logic}

\subsection{Calibration sampling}

The calibration path samples a small multiplicity of gammas and injects them
with controlled energy, position, and direction. This path is used for
benchmarks, controlled studies, and debugging.

\subsection{Veto conditions}

Deposits in active xenon outside the FV veto the event when they exceed the
configured threshold. This allows the event to be rejected immediately without
continuing detailed transport.

\subsection{FV acceptance and multi-site rejection}

Once in the FV, the event is classified using the deposit pattern returned by
the full stack. The key observable is whether the deposit pattern is consistent
with a single-site or multi-site event under the configured clustering model.

\subsection{Role of \texttt{dz\_resolution}}

The FV clustering logic uses a longitudinal clustering scale
\texttt{dz\_resolution}. Deposits separated beyond this scale contribute to
multi-site structure, while deposits consistent within this scale may be merged
into a single-site candidate.

\section{Code Map}

The current code ownership is as follows.

\begin{description}[leftmargin=2.5cm,style=nextline]
\item[\texttt{geometry\_core.jl}] shared geometric primitives, wrappers, and
geometric helpers.
\item[\texttt{geometry.jl}] canonical detector geometry, tracking tree, region
semantics, fast-kernel compilation, FV geometry compilation.
\item[\texttt{tracking.jl}] fast transport outside the FV, FV-only full stack,
event-level logic.
\item[\texttt{source.jl}] source-side propagation inside source volumes.
\item[\texttt{decays.jl}] decay and calibration sampling.
\item[\texttt{scripts/test\_event\_processing.jl}] canonical benchmark and
event-processing driver.
\item[\texttt{scripts/geometry\_viewer.html}] canonical detector and source
geometry viewer.
\end{description}

The public entry points of primary interest are:

\begin{itemize}
\item \texttt{load\_tracking\_detector}
\item \texttt{compile\_fastkernel\_geometry}
\item \texttt{compile\_fv\_geometry}
\item \texttt{process\_event}
\item \texttt{process\_event\_fastkernel\_calib}
\item \texttt{propagate\_gamma\_in\_fv}
\end{itemize}

\section{Current Limitations}

The current simulation is fast and purposeful, but it is not intended to be an
arbitrary general detector framework. The main limitations are:

\begin{enumerate}[label=\arabic*.]
\item The fast kernel is specialized to the canonical detector partition.
\item The ROI-floor termination is an intentional approximation.
\item The outside-FV transport is optimized for rapid rejection rather than for
full generic cascade tracking in every material.
\item The benchmarked fast path is best suited to the current detector semantics
and event logic, not to arbitrary future geometries without recompilation of the
fast-kernel model.
\end{enumerate}

These limitations are acceptable because they follow directly from the project
goal: simulate extremely large event samples as fast as possible while retaining
the detailed physics only where it matters.

\end{document}
